\documentclass[11pt,a4paper]{article}
% Prefer XeLaTeX or LuaLaTeX; falls back safely under pdfTeX.
\usepackage{iftex}
\ifXeTeX
  \usepackage{fontspec}
  \usepackage[spanish,shorthands=off]{babel}
  \setmainfont{Latin Modern Roman}
  \IfFontExistsTF{DejaVu Sans Mono}
    {\setmonofont{DejaVu Sans Mono}}
    {\IfFontExistsTF{Consolas}{\setmonofont{Consolas}}{\setmonofont{Latin Modern Mono}}}
\else\ifLuaTeX
  \usepackage{fontspec}
  \usepackage[spanish,shorthands=off]{babel}
  \setmainfont{Latin Modern Roman}
  \IfFontExistsTF{DejaVu Sans Mono}
    {\setmonofont{DejaVu Sans Mono}}
    {\IfFontExistsTF{Consolas}{\setmonofont{Consolas}}{\setmonofont{Latin Modern Mono}}}
\else
  \usepackage[T1]{fontenc}
  \usepackage[utf8]{inputenc}
  \DeclareUnicodeCharacter{2248}{\ensuremath{\approx}}
  \usepackage{lmodern}
  \usepackage[spanish,shorthands=off]{babel}
\fi\fi

\usepackage{xcolor}
\definecolor{navyreport}{HTML}{13294B}
\usepackage[margin=23mm]{geometry}
\title{Índice general}
\author{}
\date{}
\begin{document}
\renewcommand{\labelenumii}{\theenumi.\arabic{enumii}}
\maketitle
\section*{Introducción}

El proyecto calcula rutas ciclistas sobre OpenStreetMap. La solicitud de Plaza Norte a la Universidad Nacional de Ingeniería (UNI) sirve como caso común para seguir los datos desde la entrada hasta las salidas.

El árbol muestra las carpetas de la raíz y detalla los subpaquetes y archivos seleccionados. Los elementos en \textbf{\textcolor{navyreport}{azul marino}} delimitan el alcance de la exposición. El número indica el ponente responsable de cada grupo.

\begin{center}
\footnotesize
\renewcommand{\arraystretch}{0.88}
\begin{tabular}{@{}l@{}}
\texttt{\textbf{\textcolor{navyreport}{Proyecto\_Final\_POO/}}}\\
\hspace*{0em}\texttt{\textbar-- .cache/}\\
\hspace*{0em}\texttt{\textbar-- .git/}\\
\hspace*{0em}\texttt{\textbar-- .pytest\_cache/}\\
\hspace*{0em}\texttt{\textbar-- .ruff\_cache/}\\
\hspace*{0em}\texttt{\textbar-- .venv/}\\
\hspace*{0em}\texttt{\textbar-- .vscode/}\\
\hspace*{0em}\texttt{\textbar-- build/}\\
\hspace*{0em}\texttt{\textbar-- \textbf{\textcolor{navyreport}{dist/}}} \quad (Ponente 4)\\
\hspace*{1.5em}\texttt{\textbar-- \textbf{\textcolor{navyreport}{escritorio/}}}\\
\hspace*{3em}\texttt{\textbar-- \textbf{\textcolor{navyreport}{RutasCiclistas.exe}}}\\
\hspace*{0em}\texttt{\textbar-- docs/}\\
\hspace*{0em}\texttt{\textbar-- \textbf{\textcolor{navyreport}{scripts/}}} \quad (Ponente 4)\\
\hspace*{1.5em}\texttt{\textbar-- \textbf{\textcolor{navyreport}{compilar\_windows.ps1}}}\\
\hspace*{1.5em}\texttt{\textbar-- \textbf{\textcolor{navyreport}{desktop\_entry.py}}}\\
\hspace*{0em}\texttt{\textbar-- skills/}\\
\hspace*{0em}\texttt{\textbar-- \textbf{\textcolor{navyreport}{src/}}}\\
\hspace*{1.5em}\texttt{\textbar-- \textbf{\textcolor{navyreport}{app/}}}\\
\hspace*{3em}\texttt{\textbar-- \textbf{\textcolor{navyreport}{domain/}}} \quad (Ponente 1)\\
\hspace*{4.5em}\texttt{\textbar-- \textbf{\textcolor{navyreport}{models.py}}}\\
\hspace*{3em}\texttt{\textbar-- \textbf{\textcolor{navyreport}{providers/}}}\\
\hspace*{4.5em}\texttt{\textbar-- \textbf{\textcolor{navyreport}{base.py}}}\\
\hspace*{4.5em}\texttt{\textbar-- \textbf{\textcolor{navyreport}{factory.py}}}\\
\hspace*{3em}\texttt{\textbar-- \textbf{\textcolor{navyreport}{routing/}}} \quad (Ponente 2)\\
\hspace*{4.5em}\texttt{\textbar-- \textbf{\textcolor{navyreport}{profile.py}}}\\
\hspace*{4.5em}\texttt{\textbar-- \textbf{\textcolor{navyreport}{astar.py}}}\\
\hspace*{4.5em}\texttt{\textbar-- \textbf{\textcolor{navyreport}{dijkstra.py}}}\\
\hspace*{3em}\texttt{\textbar-- \textbf{\textcolor{navyreport}{services/}}} \quad (Ponente 3)\\
\hspace*{4.5em}\texttt{\textbar-- \textbf{\textcolor{navyreport}{planner.py}}}\\
\hspace*{3em}\texttt{\textbar-- \textbf{\textcolor{navyreport}{ui/}}} \quad (Ponente 4)\\
\hspace*{4.5em}\texttt{\textbar-- \textbf{\textcolor{navyreport}{streamlit\_app.py}}}\\
\hspace*{4.5em}\texttt{\textbar-- \textbf{\textcolor{navyreport}{map\_renderer.py}}}\\
\hspace*{4.5em}\texttt{\textbar-- \textbf{\textcolor{navyreport}{cli.py}}}\\
\hspace*{3em}\texttt{\textbar-- \textbf{\textcolor{navyreport}{visualization/}}}\\
\hspace*{4.5em}\texttt{\textbar-- \textbf{\textcolor{navyreport}{plotter.py}}}\\
\hspace*{0em}\texttt{`-- tests/}\\
\end{tabular}
\end{center}

Los archivos \textbf{\textcolor{navyreport}{\texttt{nereo.vbs}}} y \textbf{\textcolor{navyreport}{\texttt{iniciar\_app.bat}}}, ubicados en la raíz, se analizan en el tema 4.

El informe sigue el orden de datos, búsqueda, integración y presentación.
\clearpage
\begin{enumerate}
\item \textbf{Datos geográficos y adquisición de la red vial} (Ponente 1)
  \begin{enumerate}
    \item Fundamentos geográficos y arquitectura de proveedores
    \item Validación de los datos de entrada
    \item Delimitación del área de descarga
    \item Adquisición de la red y selección del proveedor
  \end{enumerate}
\item \textbf{Modelado de costes y búsqueda de caminos} (Ponente 2)
  \begin{enumerate}
    \item Fundamentos de la búsqueda en grafos ponderados
    \item Implementación del perfil ciclista
    \item Exploración con A* y Dijkstra
    \item Reconstrucción del camino y cálculo de distancia
  \end{enumerate}
\item \textbf{Integración del cálculo y planificación de rutas} (Ponente 3)
  \begin{enumerate}
    \item Fundamentos de la planificación de rutas
    \item Preparación de la red y ajuste de los extremos
    \item Búsqueda, alternativas y consolidación del resultado
  \end{enumerate}
\item \textbf{Interfaces, visualización y distribución} (Ponente 4)
  \begin{enumerate}
    \item Fundamentos de ejecución y presentación
    \item Flujo común de las interfaces y sus salidas
    \item Interfaz web y visualización interactiva
    \item Interfaz de terminal y exportación PNG
    \item Construcción y distribución del ejecutable para Windows
  \end{enumerate}
\end{enumerate}

\bigskip
\noindent\textbf{Reparto en palabras simples}
\begin{itemize}
\item \textbf{Ponente 1.} Explica cómo empieza el pedido del usuario. El programa ubica el origen y el destino, determina qué zona necesita y obtiene las calles en forma de grafo.
\item \textbf{Ponente 2.} Explica qué hace el buscador con esa red. Valora los tramos, compara las opciones y encuentra un camino con A* o Dijkstra.
\item \textbf{Ponente 3.} Explica cómo se conectan el pedido y el cálculo. El planificador ajusta los puntos a la red, llama al buscador y reúne la ruta y sus métricas.
\item \textbf{Ponente 4.} Muestra cómo se usa el programa y dónde se ve el resultado. Cubre el mapa web, el resumen de terminal, el PNG y el ejecutable de Windows.
\end{itemize}
\end{document}
